\documentclass[a4paper]{article}

\usepackage[utf8]{inputenc}
\usepackage[english]{babel}
\usepackage{microtype}
\usepackage{authblk}

\usepackage[giveninits,maxbibnames=1000]{biblatex}
\bibliography{references}

\usepackage{geometry}
\geometry{a4paper,tmargin=25mm,bmargin=25mm,lmargin=30mm,rmargin=30mm}

\usepackage{baskervald}
\usepackage[T1]{fontenc}
\renewcommand*{\bibfont}{\small}

\title{An introduction to fractional calculus:\\fundamental ideas and numerics}
\author[1]{Fabio Durastante}
\affil[1]{Università di Pisa, Dipartimento di Matematica}

\date{Course 20 hrs - 10 Lectures}

\begin{document}
	
\maketitle	
	
Fractional calculus, i.e., the definition of integrodifferential operators of
non-integer order, has a long history. It runs parallel to that of modern
differential calculus as it started with the first correspondence between Leibniz, John
Bernoulli, and De l'Hôpital in which it appears as a pure formula-generating 
curiosity far removed from the geometric
interpretation of ordinary derivatives. Then, it evolved towards more modern works by Riemann,
Hardy, and Littlewood, finally reaching the contemporary efforts to extend and interpret its
objects~\cite{MillerRoss,MR0361633}. Where the possibility of doing rigorous mathematics and proving results
within the new theory may be reason enough to invest time in its study, the push
of the last few years in this endeavor is due, as often happens, to the birth of
different areas of application outside of mathematics. Consider, for example,
the modeling of the transport of fluids within highly non-homogeneous media, the
modeling of non-local phenomena, or the search for evolutionary models of a 
non-Markovian nature. These needs have led to the definition of new operators or to
the reinterpretation of classical instruments and, at the same time, to the
search for numerical methods to simulate these models efficiently and accurately.

The purpose of this course is, on the one hand, to introduce the theoretical and
analytical tools necessary to formulate the constitutive laws of these new models
(definitions of operators, existence and uniqueness theorems, regularity of the
solution, etc.), from other, to show how it is possible to construct numerical
methods for their simulation. In particular, we will focus on the definition of
fractional integral operators, with the associated fractional derivative
operations according to Riemann-Liouville, Caputo, and Riesz. With these
instruments we will first discuss the formulation of ordinary fractional
equations and their solution~\cite{DiethelmBook}, then their use in the
definition of fractional diffusion problems as a fractional partial differential
equation~\cite{PodlubnyBook}. In the final part, we will briefly touch on some questions concerning
the definition of the fractional Laplacian and some of its applications within
modeling for complex networks~\cite{ArrigoDurastante,cnaa017,RiascosBook}.

\printbibliography

\end{document}